\documentclass[./book.tex]{subfiles}

\begin{document}

    \section{指令系统}

    \begin{enumerate}
        \item ISA（指令集体系结构）规定的内容主要包括
        \begin{itemize}
            \item 指令格式、指令寻址方式、操作类型以及每种操作对应的操作数的相应规定
            \item 操作数的类型、操作数寻址方式以及是按大端方式还是小端方式存放
            \item 程序可访问的寄存器编号、个数和位数
            \item 存储空间的大小和编址方式
            \item 指令执行过程的控制方式等，包括程序计数器、条件码定义
        \end{itemize}
        \item 指令字长：一条指令所包含的二进制代码的位数，取决于操作码的长度、地址码的长度和地址码的个数。指令字长通常为字节的整数倍
        \begin{itemize}
            \item 单字长指令：指令字长等于机器字长的指令。访存1次就能将指令完整取出
            \item 半字长指令：指令字长等于半个机器字长的指令
            \item 双字长指令：指令字长等于两个机器字长的指令。需访存2次才能将指令完整取出
        \end{itemize}
        \item 扩展操作码：使操作码的长度随地址码的减少而增加，不同地址数的指令可具有不同长度的操作码
        \begin{itemize}
            \item 设地址长度为$n$，上一层留出$m$种状态，下层可扩展出$m * 2^n$种状态
            \begin{analysisbox}
                指令字长16bit，4位为基本操作码
                \begin{circlenum}
                    \item 13条三地址指令$\Rightarrow$留出$2^4 - 13 = 3$种
                    \item 40条二地址指令$\Rightarrow$留出$3 * 2^4 - 12 = 8$种
                \end{circlenum}
            \end{analysisbox}
            \item 不允许短码是长码的前缀
            \item 各指令的操作码一定不能重复
        \end{itemize}
    \end{enumerate}

    \subsection{指令的寻址方式}

    寻址方式：寻找指令或操作数有效地址的方式，即确定本条指令的数据地址及下一条待执行指令的地址的方法。分为\textbf{指令寻址}和\textbf{数据寻址}两大类

    \subsubsection{指令寻址}

    \begin{enumerate}
        \item 指令寻址：寻找下一条将要执行的指令地址
        \item 顺序寻址：通过程序计数器PC加1（{\color{red}{1条指令的长度}}），自动形成下一条指令的地址
        \item 跳跃寻址：由本条指令给出下条指令地址的计算方式。通过转移类指令实现
        \begin{itemize}
            \item 绝对转移：地址码直接指出转移目标地址
            \item 相对转移：地址码指出转移目的地址相对于当前PC值的偏移量
        \end{itemize}
    \end{enumerate}

    \subsubsection{数据寻址}

    \begin{enumerate}
        \item 数据寻址：如何在指令中表示一个操作数的地址，或怎样计算出操作数的地址
        \item 寻址特征：指明属于哪种寻址方式
        \item 形式地址$(A)$：指令中的地址码字段并不代表操作数的真实地址
        \item 有效地址$(EA)$：形式地址结合寻址方式，可以计算出操作数在存储器中的真实地址，即EA
        \item 指令 = 操作码+地址码（操作数字段），地址码中可能包含寻址特征位和形式地址
        \item 按字节编址，采用大端方式，操作数的机器数为$ABCD 00FFH$
        \begin{itemize}
            \item $ABH$对应一个字节，可以表示一个地址
            \item 操作数占$32bit = 4Byte$，即4个连续地址
            \item 大端 $\Rightarrow$ 在内存中$FFH$表示的地址等于$ABH$的地址加3
        \end{itemize}
        \item {\color{red}{内存地址是无符号数}}
        \begin{analysisbox}
            操作数采用基址寻址方式。基址寄存器的内容为$F000 0000H$，形式地址（补码表示）为$FF12H$

            内存地址无符号数，两者不能直接相加，故需要先转化为真值之后在进行加减
        \end{analysisbox}
    \end{enumerate}

    \begin{table}[htbp]
        \centering
        \renewcommand{\arraystretch}{1.5}   % 行高
        \setlength{\tabcolsep}{2pt}         % 列间距
        \small
        \begin{tabular}{|
                >{\centering\arraybackslash}m{0.10\textwidth}|
                >{\centering\arraybackslash}m{0.20\textwidth}|
                >{\centering\arraybackslash}m{0.15\textwidth}|
                >{\centering\arraybackslash}m{0.20\textwidth}|
                >{\centering\arraybackslash}m{0.20\textwidth}|
                >{\centering\arraybackslash}m{0.15\textwidth}|}
            \hline
            & \textbf{含义}                                                                     & \textbf{有效地址 }  & \textbf{优点}                                                                            & \textbf{缺点}                     & \textbf{备注}                                                                              \\
            \hline
            隐含寻址    & 不明显地给出操作数的地址，而是隐含操作数的地址                                                         &                 & 有利于缩短指令字长                                                                              & 需要增加存储操作数或隐含地址的硬件               &                                                                                          \\
            \hline
            立即（数）寻址 & 指令字中的地址字段指出的是操作数本身，也称\textbf{立即数}，采用\textbf{补码}表示                               & $A$即是操作数        & 指令在执行阶段不访存，指令执行速度最快 & $A$的位数限制了立即数的范围 & 特征：$\#$ \\
            \hline
            直接寻址    & 指令字中的形式地址$(A)$就是操作数的真实地址$(EA)$                                                  & $EA = A$        & 简单，指令在执行阶段仅需访存一次                                                                       & $A$的位数限制了该指令操作数的寻址范围，操作数的地址不易修改 & \\
            \hline
            间接寻址    & 指令的地址字段是操作数有效地址所在主存单元的地址，即操作数地址的地址                                              & $EA = (A)$      & 可扩大寻址范围，便于编制程序（用间接寻址可方便地完成子程序返回） & 指令执行阶段要多次访存（$n$次间接寻址需$n+1$次访存）& \\
            \hline
            寄存器寻址   & 指令字中的地址字段给出的是操作数所在寄存器的编号                                                        & $EA = R_i$      & 执行阶段不用访存，执行速度快；寄存器数量远小于内存单元数，所以地址码位数较少，指令字长较短 & CPU的寄存器数量有限，寻址能力有限 & \\
            \hline
            寄存器间接寻址 & 指令字中的$R_i$所指寄存器给出操作数所在主存单元的地址                                                   & $EA = (R_i)$    & 扩大了寻址范围，减少了访存次数                                                                        &                                 &                                                                                          \\
            \hline
            相对寻址    & 将PC的内容加上指令格式中的形式地址$(A)$而形成操作数的有效地址，其中$(A)$是相对于\textbf{下一条}PC值的偏移量，\textbf{补码}表示 & $EA = (PC) + A$ & 操作数的地址随PC值的变化而变化，且与指令地址之间总是相差一个固定的偏移量，因此便于程序浮动（{\color{customgreen}{一段代码在程序内部的浮动}}） && 广泛应用于\textbf{转移指令}  \\
            \hline
            基址寻址    & 将基址寄存器$(BR)$的内容加上指令字中的形式地址$(A)$而形成操作数的有效地址                                      & $EA = (BR) + A$ & 扩大寻址范围（基址寄存器的位数大于形式地址$(A)$的位数）；有利于多道程序设计，并可用于编制浮动程序（{\color{customgreen}{整段代码在内存中浮动}}） & & 面向系统，用于为多道程序或数据分配存储空间  \\
            \hline
            变址寻址    & 将变址寄存器$(IX)$的内容加上指令字中的形式地址$(A)$而形成操作数的有效地址                                      & $EA = (IX) + A$ & 扩大寻址范围；适合编制循环程序                                                                        & & 立足于用户，主要用于处理数组问题，变址寄存器的内容有用户设定。{\color{red}{地址 = $a[0]$地址 + $sizeof(<TYPE>) \times IX$}} \\
            \hline
            堆栈寻址 & & & & \\
            \hline
        \end{tabular}
    \end{table}

    \subsection{程序的机器级代码表示}

    \includegraphics[width=\linewidth]{\subfix{../../assets/408/register.png}}

    \subsubsection{常用指令}

    \begin{enumerate}
        \item X86指令集不允许两个操作数同时来自主存
        \item \textbf{\code{push}指令}：将操作数压入内存的栈，常用于函数调用
        \begin{itemize}
            \item X86默认以$4B$为单位$\Rightarrow$栈中元素固定为$32bit$
            \item 栈增长方向与内存地址增长方向相反。栈从高地址向低地址增长。即每压入一个元素，\texttt{ESP}都减小
            \item \texttt{ESP}是栈顶，入栈前先将\texttt{ESP}值减4，然后将操作数压入\texttt{ESP}指示的地址
        \end{itemize}
        \item \textbf{\code{pop}指令}：出栈。出栈前先将\texttt{ESP}指示的地址中的内容出栈，然后将\texttt{ESP}的值加4
        \item \leavevmode
        \begin{lstlisting}[language={[x86masm]Assembler}]
            sub esp, 12 # 栈顶指针 - 12
            add esp, 8 # 栈顶指针 + 8
        \end{lstlisting}
        \item \textbf{\code{imul}指令}：有符号整数乘法指令。第一个操作数必须为寄存器
        \item \textbf{\code{idiv}指令}：有符号证书除法指令，只有一个操作数，即除数，被除数是\texttt{edx:eax}中的内容（共$64bit$），操作结果有两部分：商和余数，商送到\texttt{eax}，余数则送到\texttt{edx}
        \item \textbf{\code{xor}指令}：逻辑异或操作指令（e{\color{red}{x}}clusive {\color{red}{or}}）
        \item \textbf{\code{shl}指令}：逻辑左移指令（shift left）
        \item \textbf{\code{shr}指令}：逻辑右移指令（shift right）
        \item \textbf{\code{jcondition}指令}：条件转移指令，依据CPU状态字中的一系列条件状态转移
        \begin{lstlisting}[language={[x86masm]Assembler}]
            je <label> # jump when equal
            jz <label> # jump when last result was zero
            jne <label> # jump when not equal
            jg <label> # jump when greater than
            jge <label> # jump when greater than or equal to
            jl <label> # jump when less than
            jle <label> # jump when less than or equal to
        \end{lstlisting}
        \item \textbf{\code{cmp}指令}：用于比较两个操作数的值，功能相当于\code{sub}指令，不保存操作结果，进根据运算结果设置CPU状态字中的条件码，与\code{jcondition}指令搭配使用
        \item \textbf{\code{test}指令}：对两个操作数进行逐位与运算，功能相当于\code{and}指令，不保存操作结果，进根据运算结果设置CPU状态字中的条件码，与\code{jcondition}指令搭配使用
        \item \textbf{\code{call}指令}：将下一条指令的地址（返回地址）入栈（压栈保存，在函数的栈帧顶部），然后无条件转移到由标签指示的指令，返回值放置在\texttt{EAX}中
        \item \textbf{\code{ret}指令}：实现子程序的返回机制，弹出栈中保存的指令地址，然后无条件转移到保存的指令地址
        \item \textbf{\code{enter}指令}：等价于
        \begin{lstlisting}[language={[x86masm]Assembler}]
            push ebp # 保存上一层函数的栈帧基址，即EBP旧值
            mov ebp, esp # 设置当前函数的栈帧基址，即EBP新值
        \end{lstlisting}
        \item \textbf{\code{leave}指令}：等价于
        \begin{lstlisting}[language={[x86masm]Assembler}]
            mov esp, ebp
            pop ebp
        \end{lstlisting}
        \item 栈帧
        \begin{itemize}
            \item 栈帧最底部 $\Rightarrow$ 上一层栈帧地址 $\Rightarrow$ \texttt{EBP}在上一层栈帧中的旧值
            \item \texttt{EBP}
            \item 局部变量：C语言中越靠前定义的局部变量越靠近栈顶
            \item GCC为保证数据的严格对齐而规定每个函数的栈帧大小必须是$16B$的倍数
            \item 调用参数：参数列表中越靠前的参数越靠近栈顶。可以通过\code{R[ebp] + 8}或\code{R[ebp] + 12}获取实参
            \item 栈帧最顶部 $\Rightarrow$ 返回地址
            \item \texttt{ESP}
        \end{itemize}
        \item
        \begin{analysisbox}
            \begin{lstlisting}[language={[x86masm]Assembler}]
                      804846a    39 c2    cmp dword ptr edx, eax
                      804846c    7e 0d    jle xxxxxxxx
            \end{lstlisting}
            读出的信息：
            \begin{itemize}
                \item \code{39 c2} $\Rightarrow$ 每条指令占$2B$
                \item \code{804846c - 804846a = 2} 等于第一条指令长度 $\Rightarrow$ 按字节编址
                \item \code{jle} 偏移量为$0dH$
                \item \code{jle} 采用 PC（EIP）相对寻址，若跳转，则目标地址 = 下一条指令地址 + 偏移量 = 804846C + 2 + 0DH = 804847BH
            \end{itemize}
        \end{analysisbox}
    \end{enumerate}

    \subsection{CISC \& RISC}

    \begin{enumerate}
        \item 复杂指令系统计算机（CISC，Complex Instruction Set Computer）：增强原有指令的功能，设置更为复杂的新指令实现软件功能的硬化
        \item 精简指令系统计算机（RISC，Reduce Instruction Set Computer）：减少指令种类和简化指令功能，提高指令的执行速度
    \end{enumerate}

    \begin{table}[htbp]
        \centering
        \renewcommand{\arraystretch}{1.5}   % 行高
        \setlength{\tabcolsep}{8pt}         % 列间距
        \small
        \begin{tabular}{|
                >{\centering\arraybackslash}m{2.5cm}|
                >{\centering\arraybackslash}m{7cm}|
                >{\centering\arraybackslash}m{7cm}|}
            \hline
            & \textbf{CISC}      & \textbf{RISC}         \\
            \hline
            指令系统     & 复杂、庞大              & 简单、精简                 \\
            \hline
            指令数目     & 一般大于 200 条         & 一般小于 100 条            \\
            \hline
            指令字长     & 不固定                & 定长                    \\
            \hline
            可访存指令    & 不加限制               & 只有\code{LOAD/STORE}指令 \\
            \hline
            各种指令执行时间 & 相差较大               & 绝大多数在一个周期内完成          \\
            \hline
            各种指令使用频度 & 相差较大               & 都比较常用                 \\
            \hline
            通用寄存器数量  & 较少                 & 多                     \\
            \hline
            目标代码     & 难以用优化编译生成搞笑的目标代码程序 & 采用优化的编译程序，生成代码较为高效    \\
            \hline
            控制方式     & 绝大多数以微程序控制         & 绝大多数为组合逻辑控制，效率较微程序高   \\
            \hline
            指令流水线    & 可以通过一定方式实现         & 必须实现                  \\
            \hline
        \end{tabular}
    \end{table}

\end{document}
